\begin{tikzpicture}[font=\small]
  \foreach \i/\name in {0/低通 LP, 7.6/高通 HP, 15.2/带通 BP}{
    \begin{scope}[xshift=\i cm]
      \draw[ax] (-3.2,0) -- (3.6,0);
      \node[note, anchor=north west] at (3.4,-0.05){$\omega$};
      \draw[ax] (0,-0.2) -- (0,1.95);
      \node[note, anchor=south east] at (-0.12,1.55){$|H(j\omega)|$};
      \node[anchor=south, font=\small] at (0.1,2.15){\name};
    \end{scope}
  }
  % LP
  \begin{scope}
    \fill[ssblue, opacity=0.18] (-1.7,0) rectangle (1.7,1.35);
    \draw[curve] (-3.1,0) -- (-1.7,0) -- (-1.7,1.35) -- (1.7,1.35) -- (1.7,0) -- (3.3,0);
    \node[note, anchor=north] at (1.7,-0.12){$\omega_c$};
    \node[note, anchor=north] at (-1.7,-0.12){$-\omega_c$};
  \end{scope}
  % HP
  \begin{scope}[xshift=7.6cm]
    \fill[ssteal, opacity=0.18] (-3.3,0) rectangle (-1.7,1.35);
    \fill[ssteal, opacity=0.18] (1.7,0) rectangle (3.3,1.35);
    \draw[curve3] (-3.3,1.35) -- (-1.7,1.35) -- (-1.7,0) -- (1.7,0) -- (1.7,1.35) -- (3.3,1.35);
    \node[note, anchor=north] at (1.7,-0.12){$\omega_c$};
    \node[note, anchor=north] at (-1.7,-0.12){$-\omega_c$};
  \end{scope}
  % BP
  \begin{scope}[xshift=15.2cm]
    \fill[ssred, opacity=0.18] (1.0,0) rectangle (2.7,1.35);
    \fill[ssred, opacity=0.18] (-2.7,0) rectangle (-1.0,1.35);
    \draw[curve2] (-3.3,0) -- (-2.7,0) -- (-2.7,1.35) -- (-1.0,1.35) -- (-1.0,0) --
                  (1.0,0) -- (1.0,1.35) -- (2.7,1.35) -- (2.7,0) -- (3.3,0);
    \node[note, anchor=north] at (1.0,-0.12){$\omega_1$};
    \node[note, anchor=north] at (2.7,-0.12){$\omega_2$};
  \end{scope}

  \node[note, anchor=north, align=center] at (7.6,-0.95)
    {实值信号的频率响应关于 $\omega=0$ 对称，所以每个通带都成对出现};

  % ---------- ideal LP impulse response ----------
  \begin{scope}[yshift=-4.6cm, xshift=2.2cm]
    \fill[ssred, opacity=0.10] (-3.6,-0.65) rectangle (0,1.75);
    \draw[ax] (-3.8,0) -- (4.0,0);
    \node[note, anchor=north west] at (3.8,-0.05){$t$};
    \draw[ax] (0,-0.65) -- (0,1.9);
    \node[note, anchor=south east] at (-0.12,1.5){$h(t)$};
    \draw[curve, domain=-3.6:3.8, samples=380, smooth]
      plot (\x, {1.35*sin(deg(3.6*\x+0.00001))/(3.6*\x+0.00001)});
    \node[ssred, font=\footnotesize, anchor=north] at (-1.9,-0.7){$t<0$ 就已经有响应};
  \end{scope}

  \node[note, anchor=north west, align=left] at (8.4,-3.0)
    {\textbf{理想滤波器不可实现}\\[5pt]
     $h(t)=\dfrac{\sin(\omega_ct)}{\pi t}$ 在 $t<0$ 处非零 $\Rightarrow$ \textbf{非因果}；\\[3pt]
     且它不绝对可积 $\Rightarrow$ 不 BIBO 稳定。\\[6pt]
     实际滤波器只能在「通带纹波 / 阻带衰减 / 过渡带宽度」\\
     三者之间折中去逼近这个矩形 ——\\
     巴特沃斯、切比雪夫、椭圆各自选了不同的折中点。};

  \node[note, anchor=north, align=center] at (7.6,-7.6)
    {频域相乘 $Y=HX$ $\leftrightarrow$ 时域卷积 $y=h*x$：
     滤波之所以在频域这么好用，就是因为卷积退化成了一次逐点乘法};
\end{tikzpicture}
